\documentclass[12pt, a4paper]{article}

% === PAQUETES ESENCIALES ===
\usepackage[spanish,es-noshorthands]{babel}
\usepackage[utf8]{inputenc}
\usepackage[T1]{fontenc}

\usepackage{amsmath, amssymb, amsthm}
\usepackage{geometry}
\geometry{margin=2.5cm}
\usepackage{booktabs}
\usepackage{hyperref}
\usepackage{xcolor}
\usepackage{enumitem}
\usepackage{array}
\usepackage{textcomp}
\newcommand{\grado}{$^{\circ}$}
\usepackage{longtable}

% === COLORES PERSONALIZADOS ===
\definecolor{azuloscuro}{RGB}{0,70,140}
\definecolor{naranja}{RGB}{220,100,0}
\definecolor{verdeclaro}{RGB}{0,130,80}
\definecolor{grisclaro}{RGB}{240,240,245}

% === ENTORNOS ACADÉMICOS ===
\theoremstyle{definition}
\newtheorem{definicion}{Definición}[section]
\newtheorem{teorema}{Teorema}[section]
\newtheorem{ejemplo}{Ejemplo}[section]
\newtheorem{observacion}{Observación}[section]

% === COMANDO PARA NOTAS ===
\newcommand{\nota}[1]{\medskip\noindent\colorbox{grisclaro}{\parbox{0.95\linewidth}{\textit{#1}}}\medskip}

% === DOCUMENTO ===
\title{\textbf{Óptica Geométrica}\\[0.5em]
\large Transcripción, síntesis y desarrollo matemático}
\author{[Transcripción y síntesis generada a partir de material de curso]}
\date{\today}

\begin{document}
\maketitle

\begin{abstract}
Este documento recoge, transcribe y sintetiza el contenido de un curso de Óptica Geométrica
estructurado en once grandes bloques temáticos. Se parte de los fundamentos conceptuales de la
propagación de la luz —naturaleza ondulatoria, corpuscular y cuántica— para avanzar hacia las
leyes que gobiernan la reflexión y la refracción. Se presentan los experimentos históricos
canónicos que permitieron medir la velocidad de la luz (Römer, Fizeau-Foucault, Michelson) y
los marcos teóricos que la explican (Maxwell, Einstein, Huygens-Fresnel). El documento incluye
la ley de Snell con su derivación y aplicaciones (reflexión total interna, profundidad aparente,
joya oculta, cuarto de cilindro, cilindro transparente), la construcción geométrica del rayo
refractado, y varios problemas resueltos con demostración matemática rigurosa sobre reflexión
en espejos planos. Se respeta fielmente el contenido original, añadiendo síntesis analíticas
al cierre de cada sección que conectan los conceptos entre sí y con el panorama contemporáneo
de la óptica.
\end{abstract}

\tableofcontents
\newpage

% ============================================================
\section{Óptica Geométrica: Fundamentos}
\label{sec:optica_geometrica}
% ============================================================

La luz es una radiación electromagnética de longitud de onda comprendida entre \(400\) y
\(700\,\text{nm}\). En este sentido, la propagación de la luz en los diversos medios y la
interacción con la materia debería estudiarse como una particularización de las ondas
electromagnéticas.

No obstante, la propagación se rige por leyes geométricas simples que, en general, no dependen
de la longitud de onda ni de la intensidad de radiación.

\subsection{El Concepto de Rayo}

Para la Óptica Geométrica, la luz se propaga mediante el concepto de \textit{rayos} en línea
recta en medios uniformes. Estos rayos están relacionados exclusivamente por las \textbf{leyes
de reflexión y refracción (Snell)}.

\nota{``Para que la óptica geométrica sea válida, las dimensiones de los objetos deben ser
mucho mayores que la longitud de onda de la luz.''}

\textit{\textbf{Síntesis.}} La óptica geométrica es una aproximación de largo alcance de la
óptica ondulatoria: válida cuando \(\lambda \ll d\) (dimensión característica del obstáculo o
apertura). En este régimen, la difracción es despreciable y los rayos se comportan como líneas
geométricas. Esta simplificación es sorprendentemente poderosa: explica espejos, lentes,
instrumentos ópticos y gran parte de la óptica cotidiana, sin necesidad de recurrir a la
naturaleza ondulatoria de la luz.

% ============================================================
\section{Naturaleza de la Luz}
\label{sec:naturaleza_luz}
% ============================================================

La luz presenta una naturaleza compleja y complementaria: actúa como \textbf{onda} o como
\textbf{partícula} dependiendo de cómo la observemos. A nivel histórico se formularon dos
teorías iniciales:

\subsection{Teoría Corpuscular}

\textbf{Isaac Newton (1642--1727)} postuló que las fuentes emiten pequeñas partículas que
viajan en línea recta. Esto explicaba la reflexión y el viaje en el vacío, pero fallaba en
explicar fenómenos como la interferencia o la difracción.

\subsection{Teoría Ondulatoria}

\textbf{Christiaan Huygens (1629--1695)} postuló:

\begin{enumerate}[label=\arabic*.]
  \item La luz se debe a vibraciones periódicas.
  \item La luz monocromática está formada por vibraciones sinusoidales.
  \item En el vacío viaja a rapidez constante \(c \approx 3 \times 10^8\,\text{m/s}\).
  \item Se caracteriza por \(\lambda = c/f\).
  \item Su intensidad es proporcional al cuadrado de la amplitud.
\end{enumerate}

\nota{Conclusión: La luz posee una naturaleza dual. Se comporta como onda o partícula,
según el experimento.}

\textit{\textbf{Síntesis.}} La dualidad onda-partícula no es una contradicción lógica sino
una manifestación de la mecánica cuántica: la luz no es ni onda clásica ni partícula clásica,
sino un campo cuántico cuyas excitaciones (fotones) exhiben ambos comportamientos según la
observación. La teoría corpuscular de Newton explicaba bien la reflexión especular y la
propagación rectilínea, mientras que la teoría ondulatoria de Huygens predecía correctamente
la interferencia y la difracción. La síntesis definitiva llegó con la electrodinámica cuántica
(QED) de Feynman, Schwinger y Tomonaga en el siglo XX.

% ============================================================
\section{Teoría Electromagnética}
\label{sec:teoria_em}
% ============================================================

Tras aceptar que la luz es una onda, quedaba sin esclarecer: ¿qué es lo que vibra? Se supuso
la existencia del ``éter'', pero traía enormes contradicciones.

En 1864, \textbf{James Clerk Maxwell} reemplazó esa teoría mecánica por la electromagnética:
la luz es propagación de campos eléctricos y magnéticos perpendiculares entre sí y a la
dirección de propagación.

\begin{definicion}[Espectro Electromagnético]
La luz visible es solo una pequeñísima parte de las ondas electromagnéticas, las cuales NO
necesitan soporte material para viajar (pueden propagarse en el vacío).
\end{definicion}

El espectro electromagnético abarca desde ondas de radio (\(\lambda \sim 10^2\,\text{m}\))
hasta rayos gamma (\(\lambda \sim 10^{-12}\,\text{m}\)), pasando por microondas, infrarrojo,
visible (\(400\text{--}700\,\text{nm}\)), ultravioleta y rayos X.

Años después, \textbf{Heinrich Hertz} verificó la existencia de radioondas y demostró que la
luz difería en su frecuencia. Sin embargo, los mismos experimentos mostraron el
\textbf{efecto fotoeléctrico}, imposible de explicar mediante la teoría ondulatoria.

\textit{\textbf{Síntesis.}} Las ecuaciones de Maxwell unificaron electricidad, magnetismo y
óptica en un único formalismo, prediciendo que la velocidad de la luz en el vacío es
\(c = 1/\sqrt{\varepsilon_0 \mu_0}\). Esta predicción fue confirmada experimentalmente con
gran precisión. El valor exacto adoptado hoy es \(c = 299\,792\,458\,\text{m/s}\), definido
por convención desde 1983.

% ============================================================
\section{La Naturaleza Cuántica}
\label{sec:naturaleza_cuantica}
% ============================================================

Para el efecto fotoeléctrico, la energía de los electrones emitidos no dependía de la
intensidad de la luz, sino de su \textit{frecuencia}. \textbf{Albert Einstein (1905)} lo
resolvió basándose en la hipótesis de Max Planck.

\subsection{Fotones o Cuantos de Energía}

La luz está formada por un haz de pequeños corpúsculos llamados \textbf{fotones}. A diferencia
del modelo de Newton, estos fotones concentran la energía de la onda:

\begin{equation}
  E = h \cdot f
  \label{eq:energia_foton}
\end{equation}

donde \(h = 6.626 \times 10^{-34}\,\text{J}\cdot\text{s}\) es la constante de Planck.

Este descubrimiento confirma la naturaleza \textbf{dual} onda-partícula de la luz, y rechaza
definitivamente la existencia del éter luminífero.

\nota{¿Cómo medir \(c\)? El primer intento fue de Galileo con linternas desde colinas. Su
intento falló: el limitante principal era el tiempo de reacción frente a la enorme rapidez
de la luz.}

\textit{\textbf{Síntesis.}} La ecuación \(E = hf\) es el puente entre la descripción
ondulatoria (frecuencia \(f\)) y la corpuscular (energía cuantizada del fotón). Un fotón de
luz visible con \(\lambda = 500\,\text{nm}\) tiene una energía de apenas
\(\sim 4 \times 10^{-19}\,\text{J}\), lo que ilustra por qué los efectos cuánticos de la luz
son imperceptibles en la escala macroscópica, pero dominan en la interacción con la materia a
escala atómica.

% ============================================================
\section{El Experimento de Römer}
\label{sec:romer}
% ============================================================

En 1676, el astrónomo danés \textbf{Olaus Römer}, trabajando en París bajo el mando de
Cassini, realizaba observaciones detalladas de los eclipses de Ío, la luna magmática galileana
de Júpiter, de gran actividad geométrica y una perfecta constancia orbital de
\(\approx 42.5\,\text{h}\).

\subsection{El Retraso Misterioso}

Römer notó algo desconcertante: cuando la Tierra se movía alejándose de Júpiter, los eclipses
de Ío sufrían un constante retraso. Cuando la Tierra estaba en su punto más lejano, los
eclipses se diferían hasta \textbf{22 minutos}.

Römer concluyó que el retraso era simplemente el tiempo que tardaba la luz en atravesar la
distancia adicional del diámetro de la órbita de la Tierra (\(\approx 300\,000\,000\,\text{km}\)):

\begin{equation}
  c \approx \frac{300\,000\,000\,\text{km}}{22 \times 60\,\text{s}} \approx 227\,000\,\text{km/s}
  \label{eq:c_romer}
\end{equation}

Cassini rechazó la idea (creían que \(c = \infty\)). Tuvieron que pasar años para que James
Bradley corroborara este límite universal mediante un mecanismo distinto: la aberración estelar.

\nota{La simulación espacial asociada muestra cómo la luz que viaja de Ío a la Tierra tiene
que recorrer mayor distancia cuando la Tierra está en B comparado con A. Ese recorrido
``extra'' produce el retraso apreciado por Römer.}

\textit{\textbf{Síntesis.}} Aunque el valor de Römer difería en un \(\sim 25\%\) del valor
actual (por las imprecisiones en la distancia Tierra-Sol conocida en su época), su método fue
revolucionario: demostró por primera vez que la luz tiene velocidad finita, contradiciendo la
creencia dominante. El método de Römer es un ejemplo paradigmático del uso de la astronomía
como laboratorio de física fundamental.

% ============================================================
\section{El Experimento de Fizeau-Foucault}
\label{sec:fizeau}
% ============================================================

En 1849, el físico francés \textbf{Hippolyte Fizeau} logró medir la velocidad de la luz
mediante una experiencia hecha en la Tierra. Hizo pasar un haz luminoso entre dos dientes
consecutivos de una rueda dentada. Este haz se reflejaba en un espejo situado a cierta
distancia y volvía siguiendo la misma trayectoria. La rueda giraba con velocidad angular
variable.

Cuando el rayo reflejado pasaba precisamente a través de la abertura siguiente, el tiempo que
tarda la rueda en girar el ángulo comprendido entre dos huecos sucesivos (un diente), es igual
al tiempo empleado por la luz en recorrer la distancia hasta el espejo y volver.

Obtuvo un valor de \(3.13 \times 10^8\,\text{m/s}\), ligeramente superior al valor aceptado en
la actualidad que para el vacío es \(3 \times 10^8\,\text{m/s}\).

En el mismo año, su colaborador \textbf{Léon Foucault} mejoró el método, sustituyendo la rueda
por un \textit{espejo giratorio}. Además, midió la velocidad de la luz en el agua, confirmando
que la luz es \textbf{menor en el agua que en el aire}. Esto supuso el reconocimiento de la
teoría ondulatoria frente a la corpuscular de Newton, que predecía lo contrario.

\textit{\textbf{Síntesis.}} El experimento de Fizeau fue el primero en medir \(c\) sin recurrir
a distancias astronómicas. La confirmación de Foucault de que \(c_{\text{agua}} < c_{\text{aire}}\)
fue un golpe decisivo contra la mecánica newtoniana de la luz: en el modelo corpuscular,
los corpúsculos debían acelerarse al entrar en un medio más denso (para explicar la refracción),
prediciendo velocidades mayores en el agua, exactamente lo opuesto a lo observado.

% ============================================================
\section{El Experimento de Michelson}
\label{sec:michelson}
% ============================================================

Pocas mediciones llevadas a cabo por el hombre han sido tan precisas como la medida de la
velocidad de la luz. El premio Nobel se dio en 1907 a \textbf{Albert Michelson}, quien realizó
una medición extremadamente cuidadosa mejorando el método de Foucault con mejores instrumentos
y mayor tamaño. Colocó un espejo de ocho caras (octágono).

\subsection{Principio del Método}

\textbf{Albert A. Michelson (1879)} utilizó un espejo octagonal giratorio. Un haz de luz se
refleja en la cara A del octágono, viaja \(35.4\,\text{km}\) hasta un espejo lejano y retorna
a la cara G. Si el octágono gira exactamente \(\mathbf{1/8}\) de vuelta durante el viaje de ida
y vuelta de la luz, la cara G estará perfectamente alineada para reflejar el haz verticalmente
hacia el telescopio L. Midiendo la velocidad de rotación necesaria, se calcula \(c\).

La luz de una fuente intensa se refleja en el lado A del espejo octagonal hacia un espejo curvo
a varios kilómetros. El espejo lejano la regresa a la superficie G. Dado que el octágono está
girando, solo un impulso de luz breve incide a la vez. Si el espejo da vueltas con la velocidad
apropiada, el impulso reflejado regresa cuando la superficie F está exactamente en el lugar
apropiado para reflejarlo hacia el observador. Cuando se cumple esta condición, el tiempo
necesario para recorrer la distancia total es 1/8 del tiempo requerido para dar una vuelta
entera. La velocidad \(c\) puede calcularse a partir de esta rapidez de giro.

\textit{\textbf{Síntesis.}} Michelson obtuvo \(c = 299\,796\,\text{km/s}\), con una precisión
sin precedentes para la época. Sus experimentos sucesivos (incluyendo el famoso experimento
Michelson-Morley de 1887, diseñado para detectar el éter) no solo midieron \(c\) con extrema
exactitud, sino que al no detectar variación alguna de \(c\) con la dirección del movimiento
terrestre, sentaron las bases experimentales de la Relatividad Especial de Einstein (1905).

% ============================================================
\section{Las Leyes de Snell}
\label{sec:snell}
% ============================================================

La ley de Snell (para la refracción) es una fórmula para calcular el ángulo de refracción de
la luz al atravesar la separación entre dos medios con índice de refracción distinto. El nombre
proviene del matemático neerlandés Willebrord Snel van Royen (1580--1626). Aunque formulada
para explicar luz, se aplica a todo tipo de ondas.

\subsection{Enunciado de las Leyes}

\begin{itemize}
  \item El rayo incidente, el rayo reflejado, el rayo refractado y la normal se encuentran en
        un mismo plano, llamado ``plano de incidencia''.
  \item El ángulo de incidencia es igual al ángulo de reflexión (\(i = r\)).
  \item La razón del seno del ángulo de incidencia al seno del ángulo de refracción es una
        constante igual al cociente de las velocidades de ambos medios:
\end{itemize}

\begin{equation}
  \frac{\sin\theta_1}{\sin\theta_2} = \frac{v_1}{v_2} = \frac{n_2}{n_1}
  \label{eq:snell}
\end{equation}

o equivalentemente:
\begin{equation}
  n_1 \sin\theta_1 = n_2 \sin\theta_2
  \label{eq:snell2}
\end{equation}

\subsection{Índices de Refracción}

El índice de refracción \(n\) de un medio se define como la razón entre la velocidad de la luz
en el vacío \(c\) y la velocidad de la luz en dicho medio \(v\):

\begin{equation}
  n = \frac{c}{v}
  \label{eq:indice_refraccion}
\end{equation}

Varía según la naturaleza del medio material y es una propiedad óptica intrínseca de gran
importancia. A continuación se presentan algunos valores representativos:

\begin{center}
\begin{tabular}{@{}lc@{}}
\toprule
\textbf{Material} & \textbf{Índice de refracción} \\
\midrule
Vacío             & 1 \\
Aire (*)          & 1.0002926 \\
Helio (*)         & 1.000036 \\
Hidrógeno (*)     & 1.000132 \\
Óxido de carbono (IV) & 1.00045 \\
Agua (a 20\,°C)   & 1.333 \\
Hielo             & 1.309 \\
Diamante          & 2.417 \\
Acetona           & 1.36 \\
Cloroformo        & 1.48 \\
Alcohol etílico   & 1.361 \\
Acetaldehído      & 1.35 \\
\bottomrule
\end{tabular}
\end{center}

{\small (*) Medido en condiciones normales de presión y temperatura.}

\subsection{Casos de Refracción}

La ley de Snell \(n_1 \sin\theta_1 = n_2 \sin\theta_2\) permite analizar tres casos fundamentales:

\subsubsection*{Caso 1: De medio rápido a medio lento (\(n_1 < n_2\))}

Cuando la luz pasa de un medio con \textbf{menor} índice de refracción (más rápido, como el
aire) a uno con \textbf{mayor} índice (más lento, como el agua o el vidrio), el rayo refractado
\textbf{se acerca a la normal}: \(r' < i\). Esto ocurre porque \(n_2 > n_1\) implica
\(\sin(r') < \sin(i)\), y por tanto \(r' < i\). \textit{Ejemplo}: un rayo de luz que pasa del
aire (\(n \approx 1\)) al vidrio (\(n \approx 1.5\)).

\subsubsection*{Caso 2: De medio lento a medio rápido (\(n_1 > n_2\))}

Inversamente, cuando la luz viaja de un medio más denso (mayor \(n\)) hacia uno menos denso
(menor \(n\)), recobra velocidad y el rayo \textbf{se aleja de la normal}, obteniéndose
\(r' > i\). \textit{Ejemplo}: un rayo de luz que sale del agua (\(n \approx 1.33\)) hacia el
aire (\(n \approx 1\)). Este caso es crucial porque puede generar la reflexión total interna.

\subsubsection*{Caso 3: Incidencia normal (\(i = 0^{\circ}\))}

Si el rayo incide perpendicularmente a la superficie (\(i = 0^{\circ}\)), entonces \(\sin(0) = 0\),
y por tanto \(r' = 0^{\circ}\). El rayo atraviesa la interfaz \textbf{sin desviarse}, aunque sí cambia
su velocidad de propagación.

% ============================================================
\section{Reflexión Total Interna}
\label{sec:reflexion_total}
% ============================================================

Ocurre exclusivamente en el Caso 2 (\(n_1 > n_2\)): cuando el ángulo de incidencia alcanza un
valor tal que el rayo refractado rasaría la superficie (\(r' = 90^{\circ}\)), se define el
\textbf{ángulo crítico} \(i_c\). Para ángulos \(i > i_c\), no hay rayo refractado y toda la
luz se refleja internamente:

\begin{equation}
  \sin(i_c) = \frac{n_2}{n_1}
  \label{eq:angulo_critico}
\end{equation}

\nota{Ejemplo: agua (\(n \approx 4/3\)) hacia aire (\(n \approx 1\)):
\(i_c = \sin^{-1}(3/4) \approx 48.38^{\circ}\). Esto explica por qué un pez ve un ``cono de luz''
al mirar hacia arriba.}

\textit{\textbf{Síntesis.}} La reflexión total interna es la base de las fibras ópticas
modernas: la luz se propaga a lo largo de un núcleo de vidrio o plástico de alto índice
rodeado de un revestimiento de índice menor, rebotando sin pérdidas a lo largo de kilómetros.
Esta tecnología sustenta la infraestructura de internet de alta velocidad y las comunicaciones
transcontinentales actuales.

% ============================================================
\section{Principio de Huygens-Fresnel}
\label{sec:huygens}
% ============================================================

Es un método de análisis aplicado a problemas de propagación de ondas, llamado así en honor a
los físicos \textbf{Christiaan Huygens} y \textbf{Augustin-Jean Fresnel}.

\begin{definicion}[Principio de Huygens]
Todo punto de un frente de onda inicial puede considerarse como una \textbf{fuente de ondas
esféricas secundarias} que se expanden en todas las direcciones con la misma velocidad,
frecuencia y longitud de onda que el frente de onda del que proceden.
\end{definicion}

Esto permite explicar fenómenos como la reflexión, la refracción y la difracción. Por ejemplo:
si dos cuartos se conectan con una puerta y se produce sonido, la persona al otro lado lo
escuchará como si la puerta misma fuera la fuente.

\subsection{Historia y Difracción}

En 1678, Huygens propuso que cada punto alcanzado se vuelve fuente. Supuso que las secundarias
viajan ``hacia adelante'', dando una explicación cualitativa de la propagación esférica lineal.
Explicó refracción pero no pudo explicar por qué la luz se desvía al esquivar obstáculos.
Posteriormente \textbf{Fresnel} aportaría el concepto definitivo para justificar estos patrones
asombrosos: la \textbf{Difracción}.

\subsection{Difracción y Simulación}

La interferencia entre infinidades de frentes de ondas secundarias (distantes del frente central
móvil) origina los máximos luminosos y mínimos (sombra) observables como franjas de difracción
al pasar por rendijas estrechas. El mismo fenómeno ocurre cuando la luz pasa el borde de un
obstáculo, aunque no es fácilmente detectable a simple vista debido a la corta longitud de onda
de la luz visible frente al tamaño cotidiano de los objetos. El experimento de la \textbf{doble
rendija} para evidenciar este comportamiento netamente ondulatorio.

\subsubsection*{Principio de Huygens (1678)}

Christiaan Huygens (1678) propuso que cada punto de un frente de onda actúa como
\textbf{fuente de ondas esféricas secundarias} (wavelets). La envolvente de estas wavelets
forma el nuevo frente de onda. Esto explica la reflexión (\(\theta_i = \theta_r\)) y la
refracción (Ley de Snell: \(n_1 \sin\theta_i = n_2 \sin\theta_r\)) al pasar de un medio a otro.

\textit{\textbf{Síntesis.}} El principio de Huygens-Fresnel establece el puente entre la
óptica geométrica (rayos) y la óptica ondulatoria (frentes de onda): los rayos son siempre
perpendiculares a los frentes de onda. La adición de la condición de interferencia por parte
de Fresnel convirtió el principio cualitativo de Huygens en una herramienta cuantitativa que
predice correctamente los patrones de difracción. Kirchhoff formalizó matemáticamente este
principio en 1883 mediante la integral de difracción de Fresnel-Kirchhoff.

% ============================================================
\section{Construcción del Rayo Refractado}
\label{sec:construccion_rayo}
% ============================================================

Un método visual y práctico para el rayo incidente \(PO\) consiste en:

\begin{enumerate}
  \item Desde \(O\) dibujar semicírculos de radios \(n_1\) y \(n_2\).
  \item Prolongar el rayo \(PO\) hasta tocar el círculo \(n_1\) en el punto \(B\).
  \item Trazar línea perpendicular a la superficie de separación desde \(B\) que corte el
        círculo \(n_2\) en \(C\).
  \item El rayo final \(OC\) es el rayo verdaderamente refractado.
\end{enumerate}

\textit{\textbf{Síntesis.}} Esta construcción geométrica es la versión visual de la ley de
Snell y tiene una elegante interpretación ondulatoria: los radios \(n_1\) y \(n_2\) representan
las velocidades de fase en cada medio (inversamente proporcionales al índice), y la
perpendicularidad garantiza la conservación de la componente tangencial del vector de onda,
que es precisamente lo que expresa la ley de Snell en términos de vectores de onda.

% ============================================================
\section{Problemas Resueltos: Reflexión en Espejos Planos}
\label{sec:espejos}
% ============================================================

\subsection{Ejemplo 3.3: Longitud Mínima del Espejo}
\label{sec:longitud_espejo}

\begin{ejemplo}[Longitud del Espejo]
Encontrar la longitud mínima \(h\) de un espejo necesaria para que una persona de altura \(H\)
vea su reflexión completa.
\end{ejemplo}

\textbf{Solución:} La figura muestra los pies \(F\), los ojos \(E\), y la parte superior de la
cabeza \(D\) de una persona. Para que la persona vea toda su altura, un rayo de luz (\(DAE\))
debe salir de la parte superior de la cabeza, reflejarse en el espejo en \(A\), y entrar en sus
ojos, mientras que otro rayo (\(FCE\)) debe salir de sus pies, reflejarse en el espejo en
\(C\), y entrar en sus ojos.

La persona verá una reflexión de toda su altura (incluyendo las imágenes virtuales de los puntos
\(D\) y \(F\)), si la longitud del espejo es \(h = AC\) por lo menos.

\subsubsection*{Demostración Geométrica}

De la geometría de la figura (esquema geométrico del problema), vemos que:

\begin{equation}
  \overline{AB} = \frac{1}{2}\,\overline{DE}
  \label{eq:AB}
\end{equation}

\begin{equation}
  \overline{BC} = \frac{1}{2}\,\overline{EF}
  \label{eq:BC}
\end{equation}

Entonces:
\begin{align}
  \overline{AC} &= \overline{AB} + \overline{BC} = \frac{1}{2}\,\overline{DE} +
                   \frac{1}{2}\,\overline{EF} = \frac{1}{2}\left(\overline{DE} +
                   \overline{EF}\right) = \frac{1}{2}\,\overline{DF}
  \label{eq:AC}
\end{align}

Con \(h = \overline{AC}\) y \(H = \overline{DF}\) se obtiene finalmente:

\begin{equation}
  \boxed{h = \frac{1}{2}\,H}
  \label{eq:h_espejo}
\end{equation}

\nota{La altura del espejo debe ser la mitad de la altura de la persona. Esto es
independiente de la distancia a la que se ubique del espejo.}

\subsection{Espejos Planos Intersectados: \(\beta = 180^{\circ} - 2\theta\)}
\label{sec:espejos_intersectados}

\begin{ejemplo}[Espejos Planos Intersectados]
Consideremos dos superficies reflejantes de dos espejos planos que se intersecan a un ángulo
\(\theta\) (\(0^{\circ} < \theta < 90^{\circ}\)). Si un rayo luminoso incide sobre el espejo horizontal,
muestre que el rayo emergente intersectará el rayo incidente a un ángulo
\(\beta = 180^{\circ} - 2\theta\).
\end{ejemplo}

\textbf{Solución:} En el triángulo \(\Delta ABC\) de la figura, la suma de los ángulos
interiores debe ser \(180^{\circ}\), por lo cual:

\begin{equation}
  \theta + \phi + \alpha = 180^{\circ}
  \label{eq:triangulo_ABC}
\end{equation}

Este resultado puede reescribirse como:

\begin{equation}
  \phi + \alpha = 180^{\circ} - \theta
  \label{eq:phi_alpha}
\end{equation}

\subsubsection*{Desarrollo Matemático}

Similarmente, en el triángulo \(\Delta BCD\), la suma de los ángulos internos también debe ser
\(180^{\circ}\). Sabiendo que los ángulos en línea recta se complementan:

\begin{equation}
  (180^{\circ} - 2\phi) + (180^{\circ} - 2\alpha) + \beta = 180^{\circ}
  \label{eq:triangulo_BCD}
\end{equation}

Despejando \(\beta\):

\begin{align}
  360^{\circ} - 2\phi - 2\alpha + \beta &= 180^{\circ} \nonumber \\
  \beta &= 2(\phi + \alpha) - 180^{\circ}
  \label{eq:beta_intermedio}
\end{align}

Sustituyendo el hallazgo anterior (\(\phi + \alpha = 180^{\circ} - \theta\)):

\begin{align}
  \beta &= 2(180^{\circ} - \theta) - 180^{\circ} \nonumber \\
  \beta &= 360^{\circ} - 2\theta - 180^{\circ}
  \label{eq:beta_desarrollo}
\end{align}

\begin{equation}
  \boxed{\beta = 180^{\circ} - 2\theta}
  \label{eq:beta_final}
\end{equation}

\nota{Este resultado muestra que el ángulo \(\beta\) entre el rayo incidente y el emergente
depende exclusivamente del ángulo \(\theta\) entre los espejos, ¡no del ángulo de incidencia
particular!}

\subsection{Rayo Láser en Espejos Perpendiculares}
\label{sec:laser_perpendicular}

\begin{ejemplo}[Rayo Láser en Espejos Perpendiculares]
Dos espejos planos se intersecan en ángulos rectos. Un rayo láser incide en el primero de ellos
en un punto situado a \(a = 11.5\,\text{cm}\) de la intersección. ¿Para qué ángulo de
incidencia en el primer espejo el rayo incidirá en el \textbf{punto medio} del segundo espejo
(que mide \(28.0\,\text{cm}\) de largo) después de reflejarse en el primer espejo?
\end{ejemplo}

\textbf{Solución:} Con relación a la figura, la ley de reflexión de la luz en el punto \(A\)
y las condiciones del problema requieren que:

\begin{quote}
El rayo debe recorrer una distancia horizontal \(a = 11.5\,\text{cm}\) y alcanzar el punto
medio del segundo espejo, es decir, a una distancia vertical \(b = 28.0/2 = 14.0\,\text{cm}\).
\end{quote}

Por trigonometría, el ángulo de incidencia \(\theta\) se obtiene de:

\begin{equation}
  \tan\theta = \frac{a}{b} = \frac{11.5\,\text{cm}}{14\,\text{cm}} \approx 0.81242857
  \label{eq:tan_theta}
\end{equation}

Aplicando la función inversa:

\begin{equation}
  \theta \approx \tan^{-1}(0.81242857)
  \label{eq:theta_inv}
\end{equation}

\begin{equation}
  \boxed{\theta \approx 39.4^{\circ}}
  \label{eq:theta_laser}
\end{equation}

\nota{Observa que al ser espejos perpendiculares (90\grado), el rayo emergente sale paralelo al rayo
incidente pero en dirección opuesta. ¡Esto es un retro-reflector!}

% ============================================================
\section{Problemas Resueltos: Refracción (Ley de Snell)}
\label{sec:problemas_snell}
% ============================================================

\subsection{Ejemplo: Profundidad Aparente de un Pez}
\label{sec:profundidad_aparente}

\begin{ejemplo}[Profundidad Aparente]
Un pez está a una profundidad \(d\) bajo el agua. Tome el índice de refracción del agua igual a
\(4/3\). Muestre que cuando el pez es visto a un ángulo de refracción \(\theta_1\), la
profundidad aparente del pez es:
\begin{equation}
  z = \frac{3\,d\,\cos\theta_1}{\sqrt{7 + 9\cos^2\theta_1}}
  \label{eq:profundidad_aparente}
\end{equation}
\end{ejemplo}

\textbf{Solución:} Con relación a la figura, el rayo que sale del pez hacia la superficie
define dos distancias radiales: \(R\) (en agua) y \(r\) (aparente en aire). La distancia
horizontal \(x\) es la misma para ambos triángulos:

\begin{equation}
  x = R\sin\theta_2 = r\sin\theta_1
  \label{eq:distancia_horizontal}
\end{equation}

O bien:

\begin{equation}
  \frac{\sin\theta_2}{\sin\theta_1} = \frac{r}{R}
  \label{eq:razon_senos}
\end{equation}

Por la ley de Snell, al pasar del agua al aire:

\begin{equation}
  \frac{\sin\theta_2}{\sin\theta_1} = \frac{n_{\text{aire}}}{n_{\text{agua}}} =
  \frac{1}{4/3} = \frac{3}{4}
  \label{eq:snell_agua_aire}
\end{equation}

Combinando estos resultados:

\begin{equation}
  \frac{r}{R} = \frac{3}{4}
  \label{eq:r_R}
\end{equation}

\subsubsection*{Desarrollo Completo}

Además, las componentes verticales nos dan la profundidad real \(d\) y la aparente \(z\):

\begin{align}
  z &= r\cos\theta_1 \qquad ; \qquad d = R\cos\theta_2
  \label{eq:componentes_verticales}
\end{align}

Dividiendo:

\begin{equation}
  \frac{z}{d} = \frac{r\cos\theta_1}{R\cos\theta_2} = \frac{3}{4}\cdot\frac{\cos\theta_1}{\cos\theta_2}
  \label{eq:z_d}
\end{equation}

Puesto que \(\dfrac{\sin\theta_1}{\sin\theta_2} = \dfrac{4}{3}\), tenemos
\(\sin\theta_2 = \dfrac{3}{4}\sin\theta_1\). Elevando al cuadrado:

\begin{equation}
  \sin^2\theta_2 = \frac{9}{16}\sin^2\theta_1
  \label{eq:sin2}
\end{equation}

Usando solamente cosenos (\(\cos^2\theta = 1 - \sin^2\theta\)):

\begin{align}
  \cos^2\theta_2 &= 1 - \frac{9}{16}\sin^2\theta_1 = 1 - \frac{9}{16}(1-\cos^2\theta_1) \nonumber \\
  \cos^2\theta_2 &= \frac{7 + 9\cos^2\theta_1}{16} \nonumber \\
  \cos\theta_2   &= \frac{\sqrt{7 + 9\cos^2\theta_1}}{4}
  \label{eq:cos_theta2}
\end{align}

Finalmente, sustituyendo:

\begin{equation}
  \frac{z}{d} = \frac{3}{4}\cdot\frac{\cos\theta_1}{\dfrac{\sqrt{7+9\cos^2\theta_1}}{4}}
              = \frac{3\cos\theta_1}{\sqrt{7+9\cos^2\theta_1}}
  \label{eq:z_d_final}
\end{equation}

\begin{equation}
  \boxed{z = \frac{3\,d\,\cos\theta_1}{\sqrt{7 + 9\cos^2\theta_1}}}
  \label{eq:z_final}
\end{equation}

\nota{El pez siempre parece estar más cerca de la superficie de lo que realmente está. Esto
explica por qué es tan difícil pescar con arpón sin experiencia.}

\subsection{Ejemplo: Joya Oculta en Piscina}
\label{sec:joya_piscina}

\begin{ejemplo}[Joya Oculta en Piscina]
Un ladrón esconde una joya preciosa, colocándola en la parte inferior de una piscina pública.
Se coloca una balsa circular sobre la superficie del agua directamente por encima y centrada
sobre la joya. La superficie del agua está en calma. La balsa, de diámetro \(d = 4.54\,\text{m}\),
previene la joya de ser vista por cualquier observador desde encima del agua, ya sea en la balsa
o en el lado de la piscina. ¿Cuál es la profundidad máxima \(h\) de la piscina para que la joya
permanezca invisible?
\end{ejemplo}

\textbf{Solución:} El índice de refracción del agua es \(n = 4/3\). Para que la joya no sea
visible es necesario que en el punto \(A\), el rayo refractado no salga de la superficie del
agua. Aplicando la ley de Snell en el punto \(A\):

\begin{equation}
  n \cdot \sin i_c = 1 \cdot \sin 90^{\circ} = 1
  \label{eq:snell_critico}
\end{equation}

\begin{equation}
  \sin i_c = \frac{1}{n}
  \label{eq:sin_ic}
\end{equation}

Pero de la geometría del triángulo joya-\(A\)-centro:

\begin{equation}
  \sin i_c = \frac{d/2}{\sqrt{h^2 + (d/2)^2}}
  \label{eq:sin_ic_geo}
\end{equation}

Luego, igualando:

\begin{equation}
  \frac{d/2}{\sqrt{h^2 + (d/2)^2}} = \frac{1}{n} = \frac{3}{4}
  \label{eq:igualando}
\end{equation}

Despejando \(h\):

\begin{align}
  \frac{d}{2} &= \frac{3}{4}\sqrt{h^2 + \frac{d^2}{4}} \nonumber \\
  \left(\frac{d}{2}\right)^2 &= \frac{9}{16}\left(h^2 + \frac{d^2}{4}\right) \nonumber \\
  \frac{d^2}{4} &= \frac{9h^2}{16} + \frac{9d^2}{64} \nonumber \\
  \frac{d^2}{4} - \frac{9d^2}{64} &= \frac{9h^2}{16} \nonumber \\
  \frac{16d^2 - 9d^2}{64} &= \frac{9h^2}{16} \nonumber \\
  \frac{7d^2}{64} &= \frac{9h^2}{16} \implies h^2 = \frac{7d^2}{36}
  \label{eq:h2}
\end{align}

\begin{equation}
  \boxed{h = \frac{\sqrt{7}\,d}{6}}
  \label{eq:h_joya}
\end{equation}

Reemplazando valores:

\begin{equation}
  h = \frac{\sqrt{7}\,(4.54\,\text{m})}{6} \approx 2.0\,\text{m}
  \label{eq:h_joya_valor}
\end{equation}

\nota{La piscina debe tener como máximo 2.0 m de profundidad. Si fuera más profunda, la joya
sería visible desde los bordes porque los rayos escaparían del cono de reflexión total.}

\subsection{Ejemplo: Refracción en Cuarto de Cilindro}
\label{sec:cuarto_cilindro}

\begin{ejemplo}[Refracción en Cuarto de Cilindro]
Un material que tiene un índice de refracción \(n\) está rodeado por vacío y tiene la forma de
un cuarto de círculo de radio \(R\). Un rayo de luz paralelo a la base del material incide desde
la izquierda a una distancia \(L\) por encima de la base y emerge desde el material a un ángulo
\(\theta\). \textbf{Determine una expresión matemática para \(\sen\theta\)}.
\end{ejemplo}

\textbf{Solución -- Primera Frontera:} En la superficie curva de entrada (punto \(A\)), la
normal \(N_1\) pasa directamente por el centro de curvatura \(O\). De la geometría de la figura,
en el triángulo \(\Delta ABC\) conformado por el rayo refractado, encontramos la siguiente
relación de ángulos, observando que la horizontal se alinea con el ángulo incidente:

\begin{equation}
  \theta_1 = i + r' \implies r' = \theta_1 - i
  \label{eq:primera_frontera}
\end{equation}

\subsubsection*{Desarrollo Analítico}

Aplicamos la ley de Snell al entrar (vacío \(n=1\) a material \(n\)):

\begin{equation}
  1\cdot\sin\theta_1 = n\cdot\sin r' = n\cdot\sin(\theta_1 - i)
  \label{eq:snell_entrada}
\end{equation}

\begin{equation}
  i = \theta_1 - \sin^{-1}\!\left(\frac{\sin\theta_1}{n}\right)
  \label{eq:angulo_i}
\end{equation}

Aplicamos ahora Snell en la salida plana (punto \(B\)), con normal horizontal:

\begin{equation}
  n\cdot\sin i = 1\cdot\sin\theta \implies \sen\theta = n\cdot\sen\!\left[\theta_1 -
  \sin^{-1}\!\left(\frac{\sen\theta_1}{n}\right)\right]
  \label{eq:snell_salida}
\end{equation}

Expandiendo usando la resta de ángulos trigonométrica y sabiendo que
\(\cos(\sin^{-1} u) = \sqrt{1-u^2}\), y sustituyendo
\(\sin\theta_1 = \dfrac{L}{R}\) y \(\cos\theta_1 = \dfrac{\sqrt{R^2-L^2}}{R}\):

\begin{align}
  \sen\theta &= n\left[\left(\frac{L}{R}\right)\sqrt{1 - \left(\frac{L}{nR}\right)^2} -
               \left(\frac{\sqrt{R^2-L^2}}{R}\right)\left(\frac{L}{nR}\right)\right] \nonumber \\
  \sen\theta &= n\left[\frac{L}{R}\sqrt{\frac{n^2R^2 - L^2}{n^2R^2}} -
               \frac{L\sqrt{R^2-L^2}}{nR^2}\right] \nonumber \\
  \sen\theta &= \frac{L}{nR^2}\left[\sqrt{n^2R^2-L^2} - \sqrt{R^2-L^2}\right]
  \label{eq:seno_theta}
\end{align}

Factorizando \(\dfrac{1}{nR^2}\) y simplificando:

\begin{equation}
  \boxed{\sen\theta = \frac{L}{R^2}\left(\sqrt{n^2R^2-L^2} - \sqrt{R^2-L^2}\right)}
  \label{eq:seno_theta_final}
\end{equation}

\nota{¡Cuidado al estudiar en otros foros! Frecuentemente en la literatura se factoriza mal
perdiendo dimensionalidad (pasan la \(L\) a dividir toda la raíz erróneamente). Esta es la
expresión matemáticamente rigurosa.}

\subsection{Ejemplo: Reflexión en Cilindro Transparente}
\label{sec:cilindro}

\begin{ejemplo}[Reflexión en Cilindro Transparente]
Un cilindro transparente con un radio \(R = 2.00\,\text{m}\) tiene una \textbf{superficie de
espejo} en su mitad derecha. Un rayo de luz que se desplaza en el aire incide sobre el lado
izquierdo del cilindro. El rayo de luz incidente y el rayo de luz saliente son paralelos y distan
\(d = 2.00\,\text{m}\). \textbf{Determine el índice de refracción} del material del cilindro.
\end{ejemplo}

\textbf{Demostración Completa:} De la figura anterior se observa que el ángulo de incidencia
\(i\) es un ángulo externo al triángulo \(\Delta ACB\). Por lo cual:

\begin{equation}
  i = 2\theta
  \label{eq:i_2theta}
\end{equation}

O bien:

\begin{equation}
  \theta = \frac{i}{2}
  \label{eq:theta_medio}
\end{equation}

Además, en la figura geométrica también se observa que la altura total \(d\) se relaciona como:

\begin{equation}
  R\cdot\sin i = \frac{d}{2}
  \label{eq:R_seni}
\end{equation}

De esta manera despejamos:

\begin{equation}
  \sin i = \frac{d}{2R}
  \label{eq:seni}
\end{equation}

Por tanto:

\begin{equation}
  i = \sin^{-1}\!\left(\frac{d}{2R}\right) = \sin^{-1}\!\left(\frac{2.00\,\text{m}}{2(2.00\,\text{m})}\right)
    = \sin^{-1}\!\left(\frac{1}{2}\right) = 30^{\circ}
  \label{eq:calculo_i}
\end{equation}

\begin{equation}
  \theta = i/2 = 15^{\circ}
  \label{eq:calculo_theta}
\end{equation}

Aplicando la \textbf{segunda ley de Snell} en el punto de incidencia \(A\) (\(n_{\text{aire}} = 1\)):

\begin{equation}
  1\cdot\sin i = n\cdot\sin\theta
  \label{eq:snell_cilindro}
\end{equation}

Despejando el índice de refracción:

\begin{equation}
  n = \frac{\sin i}{\sin\theta} = \frac{\sin(30^{\circ})}{\sin(15^{\circ})}
  \label{eq:n_cilindro}
\end{equation}

\begin{equation}
  \boxed{n \approx 1.932}
  \label{eq:n_valor}
\end{equation}

% ============================================================
\section*{Resumen de Ecuaciones}
\label{sec:tabla_ecuaciones}
% ============================================================

\begin{center}
\begin{longtable}{@{}clp{5.5cm}p{3cm}c@{}}
\toprule
\textbf{\#} & \textbf{Nombre} & \textbf{Ecuación} & \textbf{Variables} & \textbf{Sec.} \\
\midrule
\endhead
1 & Energía del fotón & \(E = h \cdot f\) &
  \(h=6.626\times10^{-34}\,\text{J·s}\), \(f\): frecuencia & \S\ref{sec:naturaleza_cuantica} \\[4pt]

2 & Velocidad de Römer & \(c \approx \dfrac{300{,}000{,}000\,\text{km}}{22\times60\,\text{s}}\) &
  \(c\): velocidad de la luz & \S\ref{sec:romer} \\[4pt]

3 & Ley de Snell & \(n_1\sin\theta_1 = n_2\sin\theta_2\) &
  \(n\): índice, \(\theta\): ángulo & \S\ref{sec:snell} \\[4pt]

4 & Índice de refracción & \(n = c/v\) &
  \(v\): velocidad en el medio & \S\ref{sec:snell} \\[4pt]

5 & Ángulo crítico & \(\sin i_c = n_2/n_1\) &
  \(i_c\): ángulo crítico & \S\ref{sec:reflexion_total} \\[4pt]

6 & Longitud espejo & \(h = H/2\) &
  \(H\): altura persona & \S\ref{sec:longitud_espejo} \\[4pt]

7 & Ángulo emergente (espejos) & \(\beta = 180^{\circ} - 2\theta\) &
  \(\theta\): ángulo entre espejos & \S\ref{sec:espejos_intersectados} \\[4pt]

8 & Ángulo láser & \(\tan\theta = a/b\) &
  \(a, b\): distancias geométricas & \S\ref{sec:laser_perpendicular} \\[4pt]

9 & Profundidad aparente & \(z = \dfrac{3d\cos\theta_1}{\sqrt{7+9\cos^2\theta_1}}\) &
  \(d\): prof. real, \(\theta_1\): ángulo refracción & \S\ref{sec:profundidad_aparente} \\[4pt]

10 & Profundidad máxima (joya) & \(h = \dfrac{\sqrt{7}\,d}{6}\) &
  \(d\): diámetro de la balsa & \S\ref{sec:joya_piscina} \\[4pt]

11 & Seno refracción (cuarto cilindro) &
  \(\displaystyle\sen\theta = \frac{L}{R^2}\!\left(\sqrt{n^2R^2-L^2}-\sqrt{R^2-L^2}\right)\) &
  \(L\): altura incidencia, \(R\): radio & \S\ref{sec:cuarto_cilindro} \\[4pt]

12 & Índice cilindro transparente & \(n = \sin i/\sin\theta\) &
  \(i=30^{\circ}\), \(\theta=15^{\circ}\) para datos dados & \S\ref{sec:cilindro} \\

\bottomrule
\end{longtable}
\end{center}

\end{document}
